\section*{Chapitre \ref{ch:fm:engineering-organisation-design} --- Conception de l'organisation de l'ingénierie}

\begin{solution}{exo:fm:engineering-organisation-design:1}
Les organisations qui conçoivent des systèmes sont contraintes de produire des conceptions qui reproduisent leurs structures de communication~; un effort de conception doit donc
être organisé selon le besoin de communication.
\end{solution}

\begin{solution}{exo:fm:engineering-organisation-design:2}
$6.2/15=0.41$ pages par semaine dans les deux conceptions des vingt services, et $4.3/15=0.29$ après la consolidation~; d'après
la \cref{prop:fm:engineering-organisation-design:load}, la réaffectation ne peut pas modifier la moyenne.
\end{solution}

\begin{solution}{exo:fm:engineering-organisation-design:3}
$0.9\times6=5.4$ heures par semaine chacun~: 13.5~\% d'une semaine de quarante heures, juste plus de la moitié du cap de 25~\%.
\end{solution}

\begin{solution}{exo:fm:engineering-organisation-design:4}
La stratégie de chaque desk dépend désormais de passerelles d'ordres, de données de marché et d'entrées de pricing détenues par d'autres équipes~: les edges qui vont des stratégies aux passerelles et aux
données de marché, et des passerelles vers les contrôles de risque, traversent toutes des frontières, et elles portent des changements fréquents.
\end{solution}

\begin{solution}{exo:fm:engineering-organisation-design:5}
Intégrée~: elle reste à 3 (le desk actions garde trois personnes). Centrale~: elle tombe à 1, et les trois passerelles dépendent chacune d'une seule personne.
Consolidée~: elle tombe à 2, pour l'unique passerelle, le contrôle de risque et le gestionnaire des données de marché.
\end{solution}

\begin{solution}{exo:fm:engineering-organisation-design:6}
Intégrée~: sur les desks, là où les modèles changent. Centrale~: entre les desks et la plateforme, en tant que propriétaires des spécifications de transmission (handoff), puisque
le \omterm{def:fm:engineering-organisation-design:coord}{coût de coordination} pèse sur eux.
\end{solution}

\begin{solution}{exo:fm:engineering-organisation-design:7}
Un \omterm{def:fm:engineering-organisation-design:owner}{facteur bus} de 2 pour la passerelle, le contrôle de risque et le gestionnaire des données de marché, et une charge de 1.0 page par semaine sur chacun des deux
ingénieurs de plateforme restants.
\end{solution}

\begin{solution}{exo:fm:engineering-organisation-design:8}
Déplacer les mêmes services d'une équipe à l'autre laisse la charge moyenne inchangée et la concentre sur de petites équipes fonctionnelles~; seule la consolidation ou
la correction des services bruyants la fait baisser.
\end{solution}

\begin{solution}{pb:fm:engineering-organisation-design:1}
\begin{enumerate}
\item Voir l'\cref{exo:fm:engineering-organisation-design:1}.
\item Voir la \cref{def:fm:engineering-organisation-design:teams}~: la proximité et la rapidité face à la duplication, et une seule copie solide face aux frontières.
\item Voir la \cref{def:fm:engineering-organisation-design:qd}~; la transmission de la recherche vers la production est le point où les modèles cassent.
\item Voir les \cref{def:fm:engineering-organisation-design:coord,def:fm:engineering-organisation-design:owner}.
\item Trois desks avec cinq services chacun, cinq services partagés, 33 dépendances, 34.5 changements et 6.2 pages par semaine.
\item Voir la \cref{prop:fm:engineering-organisation-design:load}.
\item Intégrée~: 22.5 changements par semaine, 0.43 page par ingénieur au maximum, \omterm{def:fm:engineering-organisation-design:owner}{facteur bus} de 3. Centrale~: 33.0, 0.90, 2.
\item 29.0, 0.67 et 3, avec une charge moyenne de 0.29~; une copie de chaque fonction et un ingénieur libre par desk, au prix de ce que chaque desk dépende d'une
seule équipe.
\item Au plus 25~\% du temps d'astreinte, deux incidents par poste de douze heures~; l'\omterm{def:fm:engineering-organisation-design:teams}{équipe centrale} la plus chargée consomme 13.5~\% d'une semaine, les desks intégrés environ
6~\%.
\item Voir l'\cref{exo:fm:engineering-organisation-design:5}.
\item Les services au \omterm{def:fm:engineering-organisation-design:owner}{facteur bus} minimal~: les passerelles dans la conception centrale.
\item D'après les personnes qui ont modifié, déployé ou traité des incidents sur chaque service au cours de l'année passée.
\item Intégrer les stratégies et le pricing~; consolider et centraliser les passerelles, les contrôles de risque et les données de marché~; garder centraux les services partagés.
\item Consolider une fonction à la fois, en formant d'abord son équipe et en ne retirant les anciennes copies qu'une fois la nouvelle exécutée en parallèle.
\item Fréquence de déploiement, délai de mise en production des changements, taux d'échec des changements, délai de rétablissement du service.
\item Une panne ou un départ arrête tous les desks~; maintenir son \omterm{def:fm:engineering-organisation-design:owner}{facteur bus} à trois ou plus, tester les changements avec des canaris, et disposer d'une solution de repli testée.
\item À chaque nouveau desk, nouvelle plateforme de négociation ou service majeur, et chaque année.
\item Elle ajoute des frontières et concentre la charge sans supprimer aucun service.
\item Intégrée~: 22.5 changements par semaine, 0.43 page par ingénieur et par semaine, \omterm{def:fm:engineering-organisation-design:owner}{facteur bus} de 3~; centrale~: 33.0, 0.90, 2.
\item Garder intégré le travail sur les stratégies et consolider les passerelles, les contrôles de risque et les gestionnaires de données de marché dupliqués sous une seule \omterm{def:fm:technology-strategy:teams}{équipe plateforme} d'au moins
trois~; mesurer la coordination, la charge d'astreinte et le \omterm{def:fm:engineering-organisation-design:owner}{facteur bus} avant et après chaque étape.
\end{enumerate}
\end{solution}

\begin{solution}{iq:fm:engineering-organisation-design:1}
Les systèmes reproduisent la structure de communication des organisations qui les construisent~; par exemple une passerelle d'ordres par desk là où les desks ne partagent pas
leurs ingénieurs.
\iqlookfor{la thèse et un exemple concret.}
\end{solution}

\begin{solution}{iq:fm:engineering-organisation-design:2}
Transforme les modèles en code de production, assume leur performance et leur exactitude, et comprend suffisamment le modèle pour le modifier sans risque.
\iqlookfor{la transmission (handoff).}
\end{solution}

\begin{solution}{iq:fm:engineering-organisation-design:3}
Pour chaque fichier ou service, compter les personnes qui l'ont modifié de façon significative sur une fenêtre récente~; le \omterm{def:fm:engineering-organisation-design:owner}{facteur bus} est le plus petit ensemble de personnes dont
le départ laisse une partie sans personne, approché par le décompte minimal.
\iqlookfor{connaissance fondée sur l'historique.}
\end{solution}

\begin{solution}{iq:fm:engineering-organisation-design:4}
Mesurer quels services alertent et pourquoi, corriger ou faire taire les plus bruyants, ajouter des procédures d'exploitation (runbooks) et de l'automatisation, et rééquilibrer la rotation~; deux par nuit dépassent
un budget soutenable.
\iqlookfor{corriger la source, puis la rotation.}
\end{solution}

\begin{solution}{iq:fm:engineering-organisation-design:5}
Pour ce qui rend le desk différent, oui~; pour ce dont chaque desk a besoin sous la même forme, non~: les plateformes partagées évitent trois copies.
\iqlookfor{la répartition intégrée-centrale par service.}
\end{solution}

\begin{solution}{iq:fm:engineering-organisation-design:6}
Choisir un propriétaire pour chaque service et une équipe pour chaque ingénieur afin de minimiser une somme pondérée de la coordination (les edges entre équipes pondérés par la fréquence des
changements), de la charge d'astreinte au-delà d'un budget et d'une pénalité pour les \omterm{def:fm:engineering-organisation-design:owner}{facteurs bus} sous un seuil~; un programme en nombres entiers.
\iqlookfor{les trois mesures comme objectif et contraintes.}
\end{solution}
